\songtitle{Corona de alambre}

\tag*{2024}\tag{narcocorrido}\tag{ricochets}\tag{nestor-mojica}\tag{spanish-lyrics}

\begin{notes}
Based on characters from \emph{Ricochets} by Carlos Thompson.\\
Character: Néstor Mojica, "Rey Pilos."\\
Co-written by Carlos Thompson, Claude, and Gemini. %
Highly guided and interacted.
\end{notes}

\begin{style}
Narcocorrido norteño with brisk bajo sexto, tuba, accordion fills and corridos tumbados punch; verse rides tense but steady with tight rhyming lines, pre-chorus tightens on snare rolls and shouted asides, chorus lands in a chantable minor-key hook. %
Lead vocal is rough and nasal with doubled ends, group shouts on the hook word, short accordion turns between phrases, and a dusty, bright mix with dry kick, biting brass, and a final lift into an ominous break.
\end{style}

\begin{lyrics}

\songpart{Spoken Intro}
\annotation{Music fades in with a heavy, somber guitar or dark synth bass}\\
Esta es la historia de Rey Pilos, ganadero, narco y paramilitar.\\
El hombre que se quiso legalizar y conservar la operación."

\songpart{Verse 1}
Mojica era Rey Pilos,\\
con la mira en salir.\\
Quiso entrar por la rendija\\
que el gobierno abrió al rendir.

Se sentó con Don Pablo,\\
para mover el tablero.\\
Legalizar la historia\\
y cambiarse el sombrero.

\songpart{Pre-Chorus 1}
Pero antes de dar el paso,\\
había cuentas que limpiar.\\
Los Diablos de Don Pablo\\
ya buscaban dónde golpear.

La calle pidió silencio,\\
y el negocio, precisión.\\
Quien estorba en la jugada\\
sale del mapa, sin razón.

\songpart{Chorus}
Reino de papel, corona de alambre,\\
se compra el futuro matando el hambre.\\
Rey Pilos se va con un pacto en la mano,\\
pero la ley no olvida la sangre del hermano.\\
Nadie se salva, la cuenta es exacta:\\
cuando el pasado cobra, la tumba es la que pacta.

\songpart{Verse 2}
En un parking de Miami\\
Samuel Gunther cayó;\\
entre ráfagas y balas\\
madre e hija vieron caer.\\
Héctor Mendoza llegaba,\\
la gerencia de Inter-Flota asumió,\\
pero la ruta del sureste\\
ya no se pudo defender.

Mojica quemó ese puente,\\
lo soltó sin compasión.\\
Si la ruta ya era veneno,\\
mejor tumbarla al montón.

\songpart{Pre-Chorus 2}
\annotation{Conversational}\\
Lejos de Miami, la guerra estalló,\\
y a los socios de Mojica la muerte alcanzó:\\
un camión bomba a Rodríguez le frenó el carruaje,\\
y en Colombia a Arbeláez le cobraron el peaje.

Morocho movía el plano, Don Pablo dio la señal,\\
usando al agente Bard como ficha electoral.\\
Le entregaron las sobras a la ley federal,\\
limpiando el terreno en un juego político y criminal.

\songpart{Chorus}
Reino de papel, corona de alambre,\\
se compra el futuro matando el hambre.\\
Rey Pilos se va con un pacto en la mano,\\
pero la ley no olvida la sangre del hermano.\\
Nadie se salva, la cuenta es exacta:\\
cuando el pasado cobra, la tumba es la que pacta.

\songpart{Verse 3}
Lo que no vio Rey Pilos\\
fue el archivo en la oficina:\\
el fundador de Inter-Flota\\
ya sabía la cocina.

Y un abogado de auditoría\\
halló la red del sureste.\\
Cada cobro, cada ruta,\\
cada puerta y cada oeste.

No querían quemar el tronco,\\
solo las ramas del jardín,\\
pero el informe completo\\
puso a la DEA en el festín.

\songpart{Bridge}
Se firmó con traje serio,\\
pero olía a confesión.\\
La legalidad de un reino\\
a veces cuesta la prisión.

Y en la mesa de los nombres\\
nadie compra el corazón.\\
Porque el orden que se inventa\\
también paga su traición.

\songpart{Chorus}
Reino de papel, corona de alambre,\\
se compra el futuro matando el hambre.\\
Rey Pilos se va con un pacto en la mano,\\
pero la ley no olvida la sangre del hermano.\\
Nadie se salva, la cuenta es exacta:\\
cuando el pasado cobra, la tumba es la que pacta.

\songpart{Outro, (The music slows down, leaving just a heavy drum beat)}
Dos años dieron la vuelta... %
y llegó la extradición.\\
El gobierno los mandó lejos... %
con la ley en el avión.
\annotation{Fade out}

\end{lyrics}

